\chapter{绪论}

本研究开头部分对研究背景和核心价值进行阐述。在\nolinebreak[3]“大众创业、万众创新”的国家战略下，虽然创新创业教育取得了较好的效果，但是实践中仍然存在诸多问题，急需解决的还有许多重要问题。本研究用理论分析与实证检验相结合的方式来探讨学术意义，并评价研究成果的应用前景。本研究对研究方法及技术路线进行了详细的阐述，并给出其技术路线，构建起逻辑框架，将各章的内容进行系统地安排，从而给后文的论证赋予了科学支撑，同时也为有关领域的实践探究给予了较为充分的参照与借鉴。

\section{研究背景}

在\nolinebreak[3]“大众创业、万众创新”政策持续深入推进的背景下，创新创业活动已成为推动经济社会发展的关键动力。王洪才\nolinebreak[3]（2021）指出，创新创业教育是中国特色高等教育体系的重要组成部分，在推进教育现代化、提升人才培养质量方面发挥着重要的战略作用\cite{wanghongcai2021}。高校课程体系作为创新创业教育的主阵地，已成为培育大学生创业意识的核心载体。

当前，高校创新创业教育课程呈现出多元化发展的鲜明特点。李其容等\nolinebreak[3]（2024）指出，创业激情的研究为理解创业者心理机制提供了重要视角，但教育模式对学生创业意愿的影响机制仍需深入探讨\cite{liqirong2024}。Manimala和Thomas\nolinebreak[3]（2017）认为，创业教育可分为理论导向型、实践导向型和两者结合型三种基本范式\cite{Manimala2017}。这一分类在国内高等教育实践中也得到了充分印证：部分高校依托传统课堂教学，侧重于通过讲授创业理论知识和思维方法来提升学生的认知水平；另一些高校则更注重企业参访、项目孵化、案例研讨和模拟训练等实践环节对学生能力的培养；少数高校开始尝试构建理论探究与实践操作有机融合的综合性课程架构，力求实现知识迁移与技能转化的深度融合\cite{Manimala2017}。

在此背景下，学生参与创新创业课程的动因呈现出明显分化。Gielnik等\nolinebreak[3]（2015）通过对学生创业培训的评估研究发现，行动调节在创业过程中发挥着关键作用\cite{Gielnik2015}。Gielnik等\nolinebreak[3]（2017）进一步从长期视角探讨了创业培训对激情的持续影响，发现系统化的培训能够有效提升学生的创业热情和持久投入\cite{Gielnik2017}。Bohlayer和Gielnik\nolinebreak[3]（2023）则关注到行动导向创业培训中创业自我效能感的下降现象，指出培训过程中可能存在的\nolinebreak[3]“训练困境”\cite{Bohlayer2023}。陈锦和张海锋\nolinebreak[3]（2025）通过对某医学院校学生的调查发现，自主动机能够显著提升学生对创新创业课程的学习成效，并增强其创业意愿\cite{chenjin2025}。这些研究表明，创业教育的效果并非简单的线性关系，而是受到多种心理机制的影响。

创业意愿的形成是个体心理特征与外部环境因素共同作用的结果。依据计划行为理论\cite{Ajzen1991}和社会认知理论\cite{Bandura1986}的分析框架，个体创业意向受到行为态度、主观规范和感知行为控制的共同影响。武梦超和李礼旭\nolinebreak[3]（2025）整合社会认知理论与计划行为理论，发现创业态度能够调节创业自我效能感与技能之间的关系，并进而增强创业意愿\cite{wumengchao2025}。罗佳等\nolinebreak[3]（2025）的研究表明，在数字化转型背景下，系统化的创业教育不仅能有效激发学生的创业积极性，还能从创业自我效能感的视角发挥双向作用\cite{luojia2025}。在创业教育研究领域，姜诗尧等\nolinebreak[3]（2024）从生命史理论视角探讨了创业者童年社会经济地位对突破式创新的影响\cite{jiangshiyao2024}，胡望斌等\nolinebreak[3]（2025）研究了创业型领导对员工内创业行为的影响\cite{huwangbin2025}，郭润萍等\nolinebreak[3]（2024）揭示了数字创业机会的客观化机理\cite{guorunping2024}，这些研究都揭示了创业行为背后的复杂心理机制。叶文平等\nolinebreak[3]（2017）发现，财富积累速度与创业者进取心之间存在显著关联，制度环境感知在其中发挥调节作用\cite{yewenping2017}。邴浩\nolinebreak[3]（2019）的实证研究表明，创业教育的效果与个体特征密切相关，明确特定教学模式的目标服务对象能够提升实践成效\cite{binghao2019}。

Al-Jubari等\nolinebreak[3]（2019）整合自我决定理论与计划行为理论，验证了自主、胜任和归属需求通过态度、主观规范和知觉行为控制间接影响创业意向的路径，模型解释了71\%的方差\cite{AlJubari2019}。Kazmi和Nábrádi\nolinebreak[3]（2017）提出，\nolinebreak[3]“为创业而教”比\nolinebreak[3]“关于创业而教”效果更强，创业教育通过创业自我效能和感知合意性影响创业意向\cite{Kazmi2017}。这些国际研究的发现为本研究提供了重要的理论参照和方法论基础。

综上所述，本研究聚焦特定教育情境中创新创业教育模式、自主动机和外部动机对创业意愿的影响路径，以天津市某高校创业学院学生群体为研究对象，采用理论推导与实证分析相结合的方法，阐释这一现象的内在机理与实践意义。

\section{研究意义}

\subsection{理论意义}

本研究的理论意义主要体现在以下三个方面：

第一，本研究将创新创业教育模式划分为理论型、实践型和融合型三种基本形态。这一分类突破了现有研究中将创业教育作为单一同质性变量的局限，为探究其与个体创业意愿的交互机制提供了系统的分析路径。关于创业教育在高等教育中的影响，Nabi等\nolinebreak[3]（2017）的系统综述指出，不同类型的创业教育项目对学生的影响存在显著差异\cite{Nabi2017}。Kaandorp等\nolinebreak[3]（2020）对学生创业者的初始网络过程研究发现，创业教育为学生提供了关键的网络构建机会\cite{Kaandorp2020}。Edelman等\nolinebreak[3]（2016）的研究则揭示了家庭支持对年轻创业者创业活动的重要影响\cite{Edelman2016}。

第二，本研究将自主动机作为核心中介变量，构建并验证了教育模式影响创业意愿的路径机制。郭润萍等\nolinebreak[3]（2024）从社会互动视角出发，通过多案例研究揭示了数字创业机会的客观化机理，强调了认知因素在创业过程中的中介作用\cite{guorunping2024}。胡望斌等\nolinebreak[3]（2025）对创业型领导影响员工内创业行为的研究发现，上行下效的机制需要借助心理因素的中介作用\cite{huwangbin2025}。创业自我效能感、外部环境支持与初创科技企业绩效的关系研究\cite{nankai2007}表明，创业自我效能感在环境支持与绩效之间发挥着重要的中介作用。基于社会全面参与的创业教育模式研究\cite{guanlixuebao2011}也发现，教育模式需要通过满足学生的心理需求来发挥作用。这些研究为本研究探讨自主动机的中介效应提供了理论支持。

第三，本研究引入外部动机作为调节变量，探究其在教育模式影响创业意愿过程中的边界条件。关于创业教育对学生创业机会识别能力的培养，创业教育的效果受到学生个体动机特征的调节\cite{EEP2025a}。对加纳某大学学生创业者的定性研究也发现，在被动空间中学生的创业行为受到外部激励因素的显著影响\cite{EEP2025b}。周冬梅等\nolinebreak[3]（2020）在\nolinebreak[3]《管理世界》发表的创业研究回顾与展望中指出，创业意愿研究需要更多关注情境因素的调节作用\cite{zhoudongmei2020}。

\subsection{现实意义}

本研究的现实意义主要体现在以下三个方面：

第一，本研究的结论对于高校构建和完善创新创业教育体系具有重要的参考价值，同时有助于厘清各类教学模式的应用边界。通过识别不同教育模式对自主动机影响的差异，高校可以根据自身条件和学生特点，选择或设计更为适宜的课程体系，从而提高创新创业教育的针对性和有效性。

第二，本研究基于自主动机与外部激励相互影响的理论框架，探讨了如何促进大学生创新创业意愿，并提出了减少对物质奖励依赖的具体路径。研究发现的\nolinebreak[3]“挤出效应”提示教育管理者，在设计激励机制时需要审慎权衡外部奖励可能带来的负面影响，应当优先考虑如何激发和维持学生的内在学习动机。

第三，本研究主要聚焦于天津市某高校创业学院及同类高校的创新创业课程体系建设、课程实施效果评估及实践教学改革措施等问题，具有较强的针对性和应用价值。研究结论可以为同类高校在课程改革、教学设计和激励政策制定等方面提供决策参考。

\section{研究方法与研究思路}

\subsection{研究方法}

本研究综合运用文献研究法、问卷调查法和统计分析法。

文献研究法。学界已形成关于创新创业教育、大学生创业意愿及其影响因素的较为系统的理论体系与实证研究成果。Liñán和Chen\nolinebreak[3]（2009）开发的创业意向量表\cite{Linan2009}、Vallerand等\nolinebreak[3]（1992）编制的学术动机量表\cite{Vallerand1992}，以及关于创业教育效果评估的相关研究，都为相关领域的理论深化与方法创新奠定了基础。

问卷调查法。本研究以天津市某高校创业学院为调研基地，面向参与创新创业课程学习的在校学生展开调查。研究对象涵盖理论导向型、实践驱动型和融合型三类学习者，采用标准化问卷收集样本数据。问卷设计以Liñán和Chen\nolinebreak[3]（2009）的创业意向量表\nolinebreak[3]（EIQ）为基础\cite{Linan2009}，结合Vallerand等\nolinebreak[3]（1992）的学术动机量表\nolinebreak[3]（AMS）的主要维度进行整合与改良\cite{Vallerand1992}。

统计分析法。本研究利用SPSSAU统计平台，依次开展问卷数据的信效度检验、描述性统计分析、变量间相关性检验及多元回归模型构建等工作。中介效应检验采用Bootstrap重采样法获取参数置信区间，以判断中介效应的显著性。

\subsection{技术路线图}

参考图\ref{fig:research_framework}所示的研究体系，本研究采用如下研究步骤来开展系统的分析工作，即问题界定、文献综述、理论模型搭建、问卷设计与修正、数据采集、数据分析、结论阐述。

\begin{figure}[H]
    \centering
    \begin{tikzpicture}[
        font=\zihao{-5},
        head/.style={draw, rectangle, minimum width=2.2cm, minimum height=0.5cm, font=\zihao{-5}, inner sep=2pt},
        item/.style={draw=none, fill=none, font=\zihao{-5}, inner sep=1pt},
        mainbox/.style={draw, rectangle, minimum width=2.4cm, minimum height=0.7cm, font=\zihao{-5}, inner sep=3pt, align=center},
        subbox/.style={draw, rectangle, minimum width=3.2cm, minimum height=0.55cm, font=\zihao{-5}, inner sep=2pt, align=center},
        dashgrp/.style={draw, dashed, inner sep=5pt, rounded corners=1pt},
        arr/.style={-{Stealth[length=1.8mm]}, semithick}
    ]
        % ===== 顶部三栏 =====
        \node[head] (bgH) {研究背景与问题};
        \node[item, below=0.05cm of bgH] (bg1) {双创战略};
        \node[item, below=0.01cm of bg1] (bg2) {教育模式多样化};
        \node[item, below=0.01cm of bg2] (bg3) {动机异质性};
        \node[draw, fit=(bgH)(bg1)(bg2)(bg3), inner sep=3pt] (bgBox) {};

        \node[head, right=1.0cm of bgH] (thH) {理论梳理};
        \node[item, below=0.05cm of thH] (th1) {计划行为理论};
        \node[item, below=0.01cm of th1] (th2) {自我决定理论};
        \node[item, below=0.01cm of th2] (th3) {社会认识理论};
        \node[draw, fit=(thH)(th1)(th2)(th3), inner sep=3pt] (thBox) {};

        \node[head, right=1.0cm of thH] (liH) {文献回顾};
        \node[item, below=0.05cm of liH] (li1) {教育模式};
        \node[item, below=0.01cm of li1] (li2) {创业意愿};
        \node[item, below=0.01cm of li2] (li3) {自主动机};
        \node[draw, fit=(liH)(li1)(li2)(li3), inner sep=3pt] (liBox) {};

        % ===== 主链：研究模型 =====
        \node[mainbox, below=1.2cm of thBox] (model) {研究模型\\构建伦理\\模型提出假设};
        % 变量选取（与 研究模型 平行）
        \node[mainbox, right=2.5cm of model] (vars) {变量选取};

        % ===== 实证研究 + 右侧子任务组 =====
        % 让 emp 与右侧 4 行子任务居中对齐：emp 放在 4 行的中点
        \node[subbox, below=2cm of vars] (e1) {问卷设计};
        \node[subbox, below=0.15cm of e1] (e2) {数据收集};
        \node[subbox, below=0.15cm of e2] (e3) {数据清洗};
        \node[subbox, below=0.15cm of e3] (e4) {信度与效度检验};
        \node[dashgrp, fit=(e1)(e2)(e3)(e4)] (eGrp) {};
        \node[mainbox] (emp) at (model |- eGrp) {实证研究};

        % ===== 数据分析 + 右侧子任务组（与 eGrp 之间留 0.6cm 间隙） =====
        \node[subbox, below=0.7cm of e4] (a1) {描述性统计与相关分析};
        \node[subbox, below=0.15cm of a1] (a2) {主效应与调节效应检验};
        \node[subbox, below=0.15cm of a2] (a3) {假设验证论点};
        \node[dashgrp, fit=(a1)(a2)(a3)] (aGrp) {};
        \node[mainbox] (ana) at (model |- aGrp) {数据分析};

        % ===== 结论与建议 + 右侧子任务组 =====
        \node[subbox, below=0.7cm of a3] (c1) {研究结论};
        \node[subbox, below=0.15cm of c1] (c2) {理论贡献};
        \node[subbox, below=0.15cm of c2] (c3) {实践启示};
        \node[dashgrp, fit=(c1)(c2)(c3)] (cGrp) {};
        \node[mainbox] (conc) at (model |- cGrp) {结论与建议};
        % ===== 研究与展望 =====
        \node[mainbox, below=0.9cm of conc] (out) {研究与展望};

        % ===== 箭头 =====
        \draw[arr] (bgBox.south) -- ++(0,-0.2) -| (model.north);
        \draw[arr] (thBox.south) -- (model.north);
        \draw[arr] (liBox.south) -- ++(0,-0.2) -| (model.north);

        \draw[arr] (model.east) -- (vars.west);
        \draw[arr] (vars.south) -- (eGrp.north);
        \draw[arr] (model.south) -- (emp.north);
        \draw[arr] (emp.east) -- (eGrp.west);
        \draw[arr] (emp.south) -- (ana.north);
        \draw[arr] (ana.east) -- (aGrp.west);
        \draw[arr] (ana.south) -- (conc.north);
        \draw[arr] (conc.east) -- (cGrp.west);
        \draw[arr] (conc.south) -- (out.north);
    \end{tikzpicture}
    \caption{研究技术路线图}
    \label{fig:research_framework}
\end{figure}

\section{研究内容}

本研究主要按照六章的架构进行论述。第一章为绪论，主要对研究背景、学术意义、方法论等做详细的阐述，并且对本研究的全文结构进行详细的规划。第二章对创新创业教育领域的主要问题进行文献综述，归纳创业意愿的影响因素以及外部激励机制研究现状，并且提出本研究的主要切入点。第三章形成理论基础体系，联系计划行为理论、自我决定理论和社会认知理论，提取出重要的研究命题，建构起完备的理论模型。第四章转入实证研究前期准备工作，即问卷的设计过程、样本的选取方法、数据预处理等，并且对量表的信度和效度进行了严格的检验。第五章给出数据分析的结果，使用SPSS软件进行变量的描述性统计、相关性检验、中介效应检验以及调节效应检验，来证明各个假设是否成立。第六章总结研究成果，研究贡献和应用价值进行阐述，对研究的不足之处进行评价，给出今后研究的方向和建议。

\clearpage
